We now consider agents with an inaccurate utility function and derive bounds on $\epsilon'$ such that the misaligned agent is an $\epsilon'$-optimizer. 

\begin{theorem} \label{thm:inaccurate_utility}
Let $A$ be a perfect utility maximizer with utility function $u$ and error $\epsilon$ with respect to the correct utility function $u^*$. Then it is a $\frac{2\epsilon}{1-\gamma}$-optimizer with respect to $u^*$. Moreover, this bound is tight.
\end{theorem}

In the random-error case when for any $\mae_{<t}$, we randomly choose $u(\mae_{<t})$ from the set $\{\max(0,u^*(\mae_{<t})-\epsilon), \min(1,u^*(\mae_{<t})+\epsilon)\}$, we give a simple lower bound that is only a factor $4$ away from the upper bound.
Consider an environment with only one perception $1$ and actions $\{0,1\}$ and $u^*(1|\mae_{<t}0) = 1-2\epsilon$ and $u^*(1|\mae_{<t}1) = 1$. With probability $1/4$, 
it holds that $u(1|\mae_{<t}1) = u(1|\mae_{<t}0)$, in which case the agent may take the suboptimal action $0$, thus losing $2\epsilon$ in instantaneous utility. At every step, it therefore loses $\epsilon/2$ instantaneous expected utility. In total, it then loses in expectation $\frac{\epsilon}{2(1-\gamma)}$. We have thus proved the following:

\begin{theorem}
Let $A$ be a perfect utility maximizer whose utility function $u$ is such that for any $e_t,\mae_{<t}a_t$, the value $u(e_t|\mae_{<t}a_t)$ is chosen independently and uniformly from the set $\{\max(0,u^*(e_t|\mae_{<t}a_t)-\epsilon), \min(1,u^*(e_t|\mae_{<t}a_t)+\epsilon)\}$. Then, the amount of utility lost is in expectation
\[
\frac{\epsilon}{2(1-\gamma)}
\]
\end{theorem}
